\documentclass[12pt, letterpaper]{article}
\usepackage[utf8]{inputenc}
\usepackage[T1]{fontenc}
\usepackage[spanish, es-tabla, english]{babel} 
\usepackage[margin=1.2in]{geometry}
\usepackage{amsmath, amssymb, amsfonts}
\usepackage[protrusion=true, expansion=true]{microtype} 
\usepackage{titlesec} 
\usepackage{enumitem}
\usepackage{xcolor}
\usepackage{titling}
\usepackage[flushmargin, hang, bottom]{footmisc} 

% --- Colores Académicos ---
\definecolor{MidnightBlue}{RGB}{25, 25, 112}
\definecolor{SlateGray}{RGB}{112, 128, 144}

% --- Tipografía Premium (Palatino Clone) ---
\usepackage{newpxtext}
\usepackage{newpxmath}

% --- Personalización de Títulos (titlesec) ---
\titleformat{\section}
  {\normalfont\large\bfseries\scshape\color{MidnightBlue}}
  {\thesection}{1em}{}[{\titlerule[0.5pt]\vspace{2pt}\titlerule[0.2pt]}]

\titlespacing*{\section}{0pt}{3.5ex plus 1ex minus .2ex}{2.3ex plus .2ex}

% --- Entorno para Reminiscencias ---
\newenvironment{reminiscence}{%
    \itshape\color{black!85}%
    \begin{quote}%
}{%
    \end{quote}%
}

% --- Comando para Bibliografía Elegante ---
\newcommand{\bibentry}[1]{\noindent\hangindent=2.5em\hangafter=1 \small #1 \par\vspace{0.75em}}

\begin{document}
\selectlanguage{english}

% --- Portada / Encabezado Estilo Libro ---
\begin{center}
    \vspace*{1cm}
    {\color{SlateGray}\textsc{Chapter 1}} \\
    \vspace{0.4cm}
    \hrule height 1.5pt
    \vspace{0.2cm}
    \hrule height 0.5pt
    \vspace{0.6cm}
    {\huge\bfseries\color{MidnightBlue} A Perspective} \\
    \vspace{0.6cm}
    \hrule height 0.5pt
    \vspace{0.2cm}
    \hrule height 1.5pt
    \vspace{0.8cm}
    {\Large\scshape Leonid Hurwicz} \\
    \vspace{1.5cm}
\end{center}

% --- Cuerpo del Texto ---
\begin{reminiscence}
Because this volume is dedicated to the memory of Elisha Pazner, let me start with two reminiscences. I first met Elisha sometime in 1969 or 1970 when I was visiting at Harvard. Elisha consulted with me on problems related to his dissertation. Initially, we did not find it easy to communicate, but I was impressed by his original insights and depth. 

Then, after 1971, I was back at Minnesota and he was visiting at Northwestern. By that time we shared the interest in problems of implementation. We had conversations on several occasions, including at least one visit by Elisha to Minneapolis. He was, I remember, particularly concerned about the relationship of the positive ``free rider'' results due to Groves and Ledyard to my earlier negative results concerning incentive compatibility (extended to public goods by Ledyard and Roberts [1974]). I feel indebted to Elisha both for the stimulation to pursue this issue and clarification of many essential points.
\end{reminiscence}

\medskip

Subsequently, Elisha's interests moved toward the implementation of rules satisfying criteria of fairness as well as efficiency in environments involving production, especially in the presence of a nontransferable endowment such as labor. His pathbreaking contributions in this area are discussed and references given in two essays of the present book (Varian and Thomson; Postlewaite).

All the essays in this book qualify under the rubric of \textit{normative economics}: they go beyond the \textit{positive economics} designed to explain the observed economic phenomena and aim at the development of criteria to be used in judging economic policies and systems. Normative economics, in turn, has focused on two major issues: one examining the logic and merits of the various welfare criteria in terms of which the consequences of economic actions may be judged, and the other examining---in the light of such criteria---the workings of various policies and mechanisms.

Among the essays in this volume, two---focused on welfare criteria (d'Aspremont; Varian and Thomson)---are of the first category; four---devoted to problems of implementation---are of the second category. Of the latter, one (Myerson) is devoted to Bayesian implementation, one (Muller and Satterthwaite) to dominance implementation, and two (Maskin; Postlewaite) to Nash implementation. Two essays (Kalai; Milgrom) are devoted to problems (bargaining, bidding) that are outside the present framework.

The first category of studies has led to the development of the concepts of orderings and other attributes of social welfare, social welfare functions, and performance (or: social choice) functions or correspondences (rules). Among the orderings of social states, the Pareto criterion is the best known and most frequently used. More recently, different notions of fairness, so important in Pazner's work, have come to play an important role. Among social welfare functions, those associated with the names of Bergson, Samuelson, and Arrow are of prime importance. Performance functions were perhaps first formalized and their importance stressed by Reiter and Mount in the context of informationally oriented models, and by Maskin in a game-theoretic (incentive-respecting) framework.\footnote{Among related precursor concepts one should mention that of a social choice function $C(S)$ in Arrow [1951, 1963] and the notion of the choice function $C(S, R)$ in Sen [1970]. It is important to note that, in its contemporary version, the social performance correspondence need not be derived from any maximization process, and that its domain is a class of environments ($n$-tuples of individual characteristics) rather than merely feasible sets and preference profiles. In particular, the status quo (e.g., initial endowment) may be part of the specification of the environment. The idea of a more general concept of a social choice function is suggested in Arrow ([1963], footnote 34, p. 104).}

In the earlier literature (Bergson [1938], Lange [1942], Lerner [1944], Arrow [1951], Koopmans [1957], Debreu [1959], Arrow and Hahn [1971]), the central question asked was whether certain mechanisms (especially the competitive or monopoly mechanism) generated Pareto-optimal allocations, and---if so---for what categories of economic environments. Subsequently, the question was reversed: instead of regarding the mechanism as given and seeking the class of environments for which it works well, one seeks mechanisms that will work well for a given class of environments. Initially, in part because of the nature of Hayek, Mises, Lange, Lerner debates concerning the feasibility of socialism, the emphasis was on the informational (as distinct from incentival) aspects of economic mechanisms. One question posed was whether, for ``classical''\footnote{Satisfying the assumptions of convexity, divisibility, and so forth as, for instance, in Koopmans's [1957] Propositions 4 and 5.} environments, there exist mechanisms with the same optimality properties as perfect competition but, in some sense, informationally more efficient than (or as efficient as) perfect competition. It was shown (Mount and Reiter [1974], Hurwicz [1977], Osana [1978]) that no mechanisms with a smaller message space would guarantee optimality. Analogous results for the Lindahl mechanism were obtained by Sato [1981]. A uniqueness result for the Walrasian mechanism was obtained by Jordan [1982] under the additional postulate of individual rationality.

But not everyone was satisfied with the analysis of performance of various mechanisms restricted to classical environments and to their purely informational properties. In particular, incentive properties in two types of nonclassical environments had engaged the economists' attention for a long time: technologies with increasing returns, and public goods. Indeed, in each of these areas alternative mechanisms had been suggested to supplement, or substitute for, the competitive market: for increasing returns, marginal cost pricing; for public goods, the Lindahl solution. (See Lindahl [1919], Hotelling [1938], Lange [1938], Lerner [1944].)

Each of these remedies was shown to have certain defects, and the analysis of these defects led to important theoretical developments. For our purposes, the problem of public goods serves as an excellent example.

\section*{Public Goods}

To begin with, standard theorems concerning the optimality properties of competitive equilibria assume the absence of public goods. Indeed, it is not obvious how the competitive equilibrium in an economy with public goods is to be defined. It seems clear, however, that any reasonable definition of a competitive market would not yield Pareto optimality in the presence of public goods. On the other hand, Lindahl equilibrium (this term is formally defined below) in public goods economies is well defined and, under the customary assumptions of convexity and so on, Pareto optimal (Foley [1970], Milleron [1972]). But, as pointed out by Samuelson [1954, 1955], it may be unrealistic to expect the behavior required of individuals in order that a Lindahl equilibrium prevail.

To understand the problem, let us illustrate the Lindahl equilibrium in Lindahl's own simple setting. In one scenario representing the Lindahl mechanism, there is an ``auctioneer'' proposing shares $p_i (i = 1, \dots, n)$, $\sum_{i=1}^n p_i = 1$, of the aggregate cost of a public service to be borne by the $n$ participants; in turn, the $i$th participant responds by specifying the level $y_i$ of the public service that would maximize his or her utility given $p_i$; agent $i$ would then contribute $p_i y_i$ to cover the costs of public service. Equilibrium obtains when the shares $p_i$ are so chosen that all agents desire the same level of public service, that is, $y_1 = \dots = y_n$.

Let the $i$th agent's utility function in an economy with two goods be $\mathring{u}^i(x^i, y)$ where $x^i$ is the amount of private good $X$ held by $i$ (after taxes) and $y$ is the level of public service $Y$ provided to every agent. Assume furthermore, as is often done, that $Y$ can be produced by using $X$ as the input; that the production function is one of constant returns; and that the units of measurement are chosen so that one unit of $X$ produces one unit of $Y$. If the contribution of $X$ (tax paid) by agent $i$ is denoted by $t^i$, we have $x^i = \bar{x}^i - t^i$ where $\bar{x}^i$ is the initial $X$-endowment of agent $i$. Assuming zero initial endowment of $Y$ and production efficiency feasibility, we obtain the balance requirement:
\[
y = \sum_{i=1}^{n} t^i.
\]

A \textit{Lindahl equilibrium} for the economy $\mathring{e} = (\mathring{e}^1, \dots, \mathring{e}^n)$, $\mathring{e}^i = (\mathring{u}^i, \bar{x}^i)$ may be defined as a vector
\[
(x^{*1}, \dots, x^{*n}, y^*; p^{*1}, \dots, p^{*n})
\]
such that
\begin{align*}
    \sum_{i=1}^{n} p_i^* &= 1 \\
    y^* &= \sum_{i=1}^{n} (\bar{x}^i - x^{*i});
\end{align*}
and for each $i = 1, \dots, n$,
\begin{gather*}
    p^{*i} \geqq 0, \quad (x^{*i}, y^*) \text{ is individually feasible,}\footnote{That is, $(x^{*i}, y^*)$ is in the $i$th agent's consumption set $C$, often chosen as the non-negative quadrant.} \\
    p_i^* y^* + x^{*i} = \bar{x}^i
\end{gather*}
and
\[
\mathring{u}^i(x^{*i}, y^*) \geqq \mathring{u}^i(x^i, y)
\]
for each individually feasible $(x^i, y)$ satisfying the budget equality $p_i^* y + x_i = \bar{x}^i$.

The preceding defines the \textit{Lindahl allocation correspondence} $L$ on the class $E$ of economies by
\begin{quote}
    $L(e) = \{(x^1, \dots, x^n, y) \colon$ for some $(p^1, \dots, p^n)$, the vector $(x^1, \dots, x^n, y; p^1, \dots, p^n)$ is a Lindahl equilibrium for $e \}$,
\end{quote}
for every $e$ in $E$. [Here, again, $e = (e^1, \dots, e^n)$, $e^i = (u^i, \bar{x}^i)$.]

The Lindahl scenario assumes that each agent treats the proposed shares $p_i$ parametrically (i.e., as given) and that the agent's response $y_i$ is based on the maximization of the agent's \textit{true} utility function. But, as stressed by Samuelson, agents might find it to their advantage to respond with a value of $y_i$ that does not maximize their true utility functions. That is, they might benefit by misrepresenting their preferences. Consequently, the Lindahl mechanism may fail to produce Pareto-optimal outcomes. If it does fail, is there another mechanism that would succeed where the Lindahl mechanism fails? In searching for such a mechanism, how is the problem to be formalized?

One answer is to be found in the theory of noncooperative games.\footnote{An important role in the development of public goods theory was played by the pioneering contributions of Drèze and de la Vallée Poussin [1969], Malinvaud [1969], and the subsequent dynamics-oriented literature. Unfortunately, I am unable to cover this work in the present essay.} Specifically, one may consider rules prescribing resource allocation behavior (such as Lindahl, or competitive profit maximization, or marginal cost pricing) as equivalent to direct-revelation mechanisms, giving rise to a special class of games (Hurwicz [1972]). But the terms ``mechanism'' and ``direct revelation'' must be defined.

A \textit{mechanism} is defined by endowing each of the $n$ participants with a \textit{strategy domain}, with $S^i$ denoting the domain of the $i$th agent, and the \textit{outcome function} (also called strategy outcome function or game form) denoted by $h$; this outcome function $h$ specifies the resource allocation $(x^1, \dots, x^n, y)$ resulting from any $n$-tuple of strategy choices. Thus the outcome function specifies the rules of the game; for the economy described previously, it may be written as follows:
\[
(x^1, \dots, x^n, y) = h(s^1, \dots, s^n), \quad s^i \in S^i, \quad i = 1, \dots, n.
\]
Formally, a mechanism is the ordered pair $(S, h)$, $S = S^1 \times \dots \times S^n$. The space of outcomes containing the range of $h$ will be noted by $Z$.

Now consider an economy $e = (e^1, \dots, e^n)$ where $e^i$ has the utility function $u^i$, and a mechanism $m = (S, h)$, $S = S^1 \times \dots \times S^n$. Let $\Gamma$ denote the noncooperative game $(S, \varphi_{m,e})$ where the $i$th agent's payoff function is given by $\varphi^i_{m,e}(s) = u^i(h(s))$ for all $s$ in $S$. Denote by $\nu_m(e)$ the set of Nash equilibria of the game $\Gamma$ for the economy $e$.\footnote{$(s^{*1}, \dots, s^{*n}) = s^*$ in $S$ is a \textit{Nash equilibrium} of the game $\Gamma = (S, \varphi)$, $S = S^1 \times \dots \times S^n$, if, for every $i = 1, \dots, n$, it is the case that $\varphi^i(s^*) \geqq \varphi^i(i, s^j/s^*)$ for all $s^j$ in $S^i$, where $i, s^j/s^*$ denotes the $n$-tuple $s^*$ with its $i$th entry $s^{*i}$ replaced by $s^j$.} We say that the mechanism $m = (S, h)$ \textit{Nash implements} a performance correspondence $F \colon E \to\to Z$ on $E$ if it is the case that for every $e$ in $E$, (1) the set $\nu_m(e)$ is not empty, and (2) $h(\nu_m(e)) \subseteq F(e)$. (This is sometimes called ``weak'' implementation.)

The general problem of designing a mechanism that yields Pareto-optimal outcomes in a given range $E$ of environments can now be formally stated as that of finding a mechanism implementing the Pareto correspondence $P \colon E \to\to Z$ on $E$ where $P(e)$ is the set of allocations that are Pareto optimal for $e$.

A \textit{direct-revelation mechanism} in $E = E^1 \times \dots \times E^n$ may be defined by the following property: $S^i = E^i$, where $E^i$ is the class of a priori admissible characteristics of agent $i$ ($i = 1, \dots, n$).

A \textit{natural}\footnote{The term \textit{natural} is introduced to focus attention on the postulated property of the outcome function.} \textit{direct-revelation mechanism} in $E = E^1 \times \dots \times E^n$ for the performance correspondence $F \colon E \to\to Z$ may then be defined as a direct-revelation mechanism in $E$ such that, for $s = e$ [i.e., for $(s^1, \dots, s^n) = (e^1, \dots, e^n)$], the outcome $h(s)$ is an element of $F(e)$; that is, $h(s)|_{s=e} \in F(e)$. In particular, when $F$ is a (single-valued) function and the mechanism $m$ is natural for $F$, it is the case that $h(s)|_{s=e} = F(e)$. That is to say, in a natural direct-revelation mechanism, the outcome is one that would be desirable according to the performance function $F$ if the agents were truthful. The corresponding game $\Gamma = (E, \varphi_{m,e})$ is called a \textit{natural direct-revelation game} for $F$.

In such a natural direct-revelation game for $F$ one would wish that the (unique) $n$-tuple of Nash equilibrium strategies be \textit{truthful}; the mechanism is then called \textit{straightforward}. This requirement can be written as
\[
\nu_\Gamma(e) = \{e\}.
\]
For clearly, in this case, $\Gamma$ Nash implements $F$. (It does not follow that a mechanism must be straightforward to implement $F$.) A natural direct-revelation mechanism whose Nash equilibria are truthful is called \textit{incentive compatible}.

Now it turns out (d'Aspremont and Gérard-Varet [1979b], Thm. 1; Dasgupta, Hammond, and Maskin [1979], Thm. 7.1.1) that if a mechanism is incentive compatible, its Nash equilibria are also dominance equilibria. But it is also known that, generally speaking, there do not exist mechanisms guaranteeing dominance equilibria for sufficiently broad classes of environments when $F$ is Pareto optimal (i.e., a subcorrespondence of the Pareto correspondence).\footnote{For economies with transferable utilities (of the form $u^i(x^i, y) = x + v_i(y)$) this non-existence follows (roughly) from the following results. (1) If for a given mechanism dominance equilibria exist, then there is an ``equivalent'' natural direct mechanism with truth as everyone's dominant strategy (Dasgupta, Hammond, and Maskin [1979], Thm. 4.1.1). (2) On ``smoothly connected'' domains $V^i$ of valuation functions $v_i$, any mechanism with truthful dominance equilibria is a ``Groves mechanism'' (Green and Laffont [1977, 1979]; Holmström [1979], Thm. 1). (3) A Groves mechanism is generically unbalanced (Green and Laffont [1979], Walker [1980], Thm. 1). (4) An unbalanced mechanism is not Pareto optimal.} Hence the aim of constructing incentive-compatible equilibria in the sense of the preceding definition is too ambitious. (This is true not only for public goods economies, but also for private goods pure exchange economies. See the subsequent discussion.)

The recognition of this fact can lead to several alternative compromises. One might sacrifice Pareto optimality or postulate probabilistic (Bayesian) beliefs. This latter approach was pioneered by Harsanyi [1967, 1968]; its economic applications include contributions by d'Aspremont and Gérard-Varet [1979a,b] and Arrow [1977]. (See the essay by Myerson, this volume.) But we shall confine ourselves here to the study of another compromise, sacrificing the dominance equilibrium property. We shall still ask for Pareto optimality and Nash implementability but no longer in a natural direct-revelation game. Thus there is no longer any necessary relationship between $S^i$ and $E^i$ and no requirement corresponding to $h(s)|_{s=e} \in F(e)$. (In fact, if the game is not of direct-revelation type, $s = e$ no longer makes sense.)

In this broader class of games, the Nash implementation of a Pareto-optimal performance correspondence is no longer impossible. This was shown by Groves and Ledyard who were the first to construct a mechanism (not of direct-revelation type) that yielded Pareto-optimal Nash allocations in public goods environments. However, the performance correspondence of the Groves--Ledyard mechanism was not individually rational. It became, therefore, natural to ask whether some Pareto optimal individually rational correspondence could be Nash implemented. Since the Lindahl performance is both Pareto optimal and individually rational, the question could be answered affirmatively by showing that the Lindahl correspondence is Nash implementable. As it turned out, it is not difficult to construct mechanisms that Nash implement the Lindahl correspondence (Hurwicz [1979], Walker [1981]).\footnote{The trailblazing contribution to the analogous problem in private goods economies is due to Schmeidler [1976], who constructed a mechanism Nash implementing the Walrasian correspondence. However, see the subsequent text concerning the individual feasibility requirement.}

For balanced implementation of the Lindahl correspondence when there are three or more agents, \textit{smooth} mechanisms ignoring the individual feasibility condition were constructed by Hurwicz [1979s]\footnote{A two-good economy (one public, one private) and constant returns technology.} and Walker [1981].\footnote{Extended to arbitrarily many public and private goods and more general technologies.} The former has price-quantity proposals as strategic messages; that is, $m_i = (p_i, y_i)$, and uses a circular arrangement of participants so that each agent $i$ is an auctioneer for his neighbor, say agent $i + 1$. Walker's arrangement is analogous but using a message space of smaller dimensions. Specifically, in a two-good, $n$-agent world with constant returns, the Walker strategy space is of dimension $n$ while the Hurwicz strategy space is of dimension $2n$. Clearly, the Walker mechanism minimizes the dimensionality of the message space.

The case of transferable utilities is somewhat special. In this case there is a unique interior Pareto-optimal level $\hat{y}$ of the public service, defined for differentiable utility functions by the Samuelson condition $\sum_{i=1}^n v'_i(\hat{y}) = 1$. (Uniqueness of $\hat{y}$ follows from the fact that $v'_i$ is assumed strictly decreasing.) For such economies it is possible to design incentive-compatible mechanisms that will yield a Pareto-optimal level $\hat{y}$ of the public service (Clarke [1971], Groves [1970], Groves and Loeb [1975], Green and Laffont [1979]). However, the balance condition cannot in general be satisfied, hence the resulting allocation, say $(x^1, \dots, x^n, \hat{y})$ is not, in general, Pareto optimal.

A (discontinuous) balanced mechanism with profile type strategy spaces where individual feasibility conditions \textit{are} satisfied is constructed in Hurwicz, Maskin, and Postlewaite [1980]. Subject to additional conditions on the preferences, a Constrained Lindahl correspondence is feasibly implemented.

The mechanisms just mentioned are balanced, in the sense that, for every $s$ in $S$, the balance requirement
\[
y = \sum_{i=1}^{n} t^i
\]
is satisfied.\footnote{More explicitly, write $(t^1, \dots, t^n, y) = h(s) = (T^1(s), \dots, T^n(s), Y(s))$; then the balance requirement is that $Y(s) = \sum_{i=1}^n T^i(s)$ for all $s \in S$.} But the requirement of individual feasibility [i.e., that $(x^i, y) \geqq 0$] might have been violated. (The same problem arose in Schmeidler's [1976] and Hurwicz's [1979] mechanisms Nash implementing the Walrasian correspondence.) However, as shown in Hurwicz, Maskin, and Postlewaite [1980], mechanisms analogous to those in Maskin [1977] can be constructed to Nash implement Lindahl, Walras, and various other correspondences without violating either the balance or the individual feasibility requirements.

These findings raise the more general question: Which correspondences are Nash implementable without violating the feasibility (individual and balance) requirements? To a considerable extent the question was answered in Maskin [1977] for cases where the feasible outcome set is a priori known to the designer of the mechanism. Maskin showed that (1) a Nash-implementable correspondence must be monotone;\footnote{$F$ is monotone when the following holds: if an outcome $z$ is $F$-desirable for a preference profile $\mathbf{R}$, and another profile $\mathbf{R}'$ is no less favorable to $z$ than $\mathbf{R}$ was, then $z$ is also $F$-desirable for $\mathbf{R}'$ (Maskin [1977]). (``$z$ is $F$-desirable for the environment $e$'' means that $z$ is an element of $F(e)$; ``$\mathbf{R}'$ is no less favorable to $z$ than $\mathbf{R}$ was'' means that $z R_i z'$ implies $z R'_i z'$ for all $i$ and all $z'$.)} and (2) if there are at least three agents, any correspondence that is monotone and has the ``no-veto power''\footnote{$F$ has the ``no-veto power'' property when the following holds: if an outcome $z$ is the most preferred one for at least $n - 1$ agents then $z$ is $F$-desirable (Maskin [1977]).} property is Nash implementable.

Analogous Nash-implementability conditions for cases where the feasible set is not a priori known to the designer are given in Hurwicz, Maskin, and Postlewaite [1980].\footnote{In particular, an example due to Postlewaite (Hurwicz, Maskin, and Postlewaite [1980]) shows that in environments where boundary equilibria can occur, the Walras correspondence is not monotone and must be replaced by its ``constrained'' counterpart.}

It is to be noted that the implementability results in Maskin [1977] and Hurwicz, Maskin, and Postlewaite [1980] are proved by constructive methods. Thus, for example, when the feasible set is a priori known to the designer and the correspondence to be implemented is monotone and has the ``no-veto power'' property, Maskin's proof of Nash implementability shows how to construct the outcome functions using preference profiles as the players' strategic variables. In Hurwicz, Maskin and Postlewaite, the profiles $(e^1, \dots, e^n)$, where $e^i$ involves endowments and/or production sets as well as preferences of the $i$ agent, are used as the strategic variables.

Economists are particularly interested in implementing correspondences that are Pareto optimal\footnote{That is, correspondences $F$ such that $F(e) \subseteq P(e)$ for all $e$ in $E$ where $P(e)$ is the set of Pareto-optimal outcomes in the environment $e$, and $E$ is the class of environments over which implementation is sought.} and either individually rational or envy free (i.e., fair). Let $F$ be Pareto optimal and individually rational. Somewhat surprisingly, it turns out (Hurwicz [1979b]) that if (i) $F$ is Pareto optimal, individually rational, and continuous, (ii) $E$ is sufficiently broad, and (iii) $F$ is Nash implementable over $E$, then (a) in public goods economies $F \supseteq L$, while (b) in private goods pure exchange economies, $F \supseteq W$ (where $W$ is the Walrasian correspondence).

An analogous result for private goods pure exchange economies concerning fairness is due to Thomson [1979], who showed that if (i) $F$ is Pareto optimal, envy free, and continuous, (ii) $E$ is sufficiently broad, and (iii) $F$ is Nash implementable over $E$, then $F \supseteq W_{I(w)}$. [Here $W_{I(w)}$ is the Walrasian correspondence from the equal initial endowment point $I(w)$ obtained by the redistribution of $w = (w^1, \dots, w^n)$; $w^i$ is $i$'s initial endowment vector.\footnote{That is, for $e = (e^1, \dots, e^n)$, $e^i = (u^i, w^i)$, $i = 1, \dots, n$, we have $z \in W_{I(w)}(e)$ if and only if $z$ is the Walrasian allocation in an economy with the same preferences $(u^1, \dots, u^n)$ but with each agent's initial endowment equalized to $I(w) = \sum_{i=1}^n w^i/n$.}]

\section*{Two-Agent Economies}

It is worth noting that special difficulties arise with regard to balanced Nash implementation when there are only two agents ($n = 2$). Most mechanisms referred to in the preceding sections work only for three or more persons. It is possible,\footnote{Ignoring the individual feasibility condition.} however, to Nash implement Walrasian and Lindahl correspondences when $n = 2$ by balanced (but discontinuous) outcome functions (Hurwicz [1979c], Miura [1982]\footnote{Miura points out an error in Hurwicz's [1979c] Lindahl-implementing outcome function; his modified outcome function does implement the Lindahl correspondence for two or more agents.}). That this cannot be done by \textit{smooth} functions is shown in Reichelstein [1984]. (An analogous result for the Pareto-optimal correspondence has recently been obtained by Hurwicz and Hans Weinberger.)

Implementation through profiles as strategies (as in Maskin [1977] and Hurwicz, Maskin, and Postlewaite [1980]) has the disadvantage of using huge (indeed, infinite-dimensional) strategy domains, and discontinuous outcome functions. But it should be understood that the profile approach provides not merely the implementation of a particular performance correspondence but, rather, an \textit{algorithm} for constructing outcome functions implementing a large class of performance correspondences.

When the specific correspondence to be implemented is known, it may be possible to do much better. In particular, for a pure exchange private goods economy, with three or more agents, and with endowments a priori known to the designer, Postlewaite and Wettstein [1983] have constructed a mechanism with a \textit{continuous} outcome function where each agent's strategy domain is of finite dimension (equal to $2l + 1$ where $l$ is the number of goods), a mechanism that Nash implements the Constrained Walrasian correspondence.

For \textit{smooth} balanced mechanisms\footnote{Ignoring the individual feasibility condition.} Reichelstein [1982] has found minimal dimensions of strategy spaces for Nash implementing the Walrasian correspondence in two-good economies with three or more traders.\footnote{This minimal dimension is lower than that used in smooth balanced mechanisms implementing the Walrasian correspondence in Hurwicz [1979].} (As mentioned previously, when there are only two traders, smooth implementation of the Walrasian correspondence is altogether impossible.)

\section*{Private Goods Economies}

Although the problem of designing an incentive-compatible mechanism was particularly obvious in economies with public goods, it also arises in private goods economies. Thus it was shown in Hurwicz [1972, Ch. 14] that in simple (two-person, two-good) pure exchange economies there exist no natural incentive-compatible mechanisms guaranteeing Pareto optimality and individual rationality. An analogous result in Dasgupta, Hammond, and Maskin [1979], Thm. 4.4.1, replaces the requirement of individual rationality by that of nondictatorship and broadens the class of admissible preferences so as to permit nonsmoothness and even discontinuities.\footnote{It appears that a strengthening of assumptions made in Dasgupta, Hammond, and Maskin [1979] is required when the number of agents exceeds two (Sonnenschein and Satterthwaite [1981], and a private communication from Eric Maskin).}

In [1983], Hurwicz and Walker show, without postulating individual rationality, that dominance-equilibrium mechanisms are generically non-optimal. This is a counterpart of the result in Walker [1980] for public goods economies and was conjectured by Walker in [1980].

Pure exchange economies were also studied by Satterthwaite [1976] and in Satterthwaite and Sonnenschein [1981]. It was shown that ``regular,'' ``nonbossy'' direct-revelation mechanisms that are strategy-proof and defined on an open set of utility functions are serially dictatorial. In certain economies with production this may generate nonoptimality.

It should be noted that the parallelism of the impossibility results, as between economies with and without public goods, breaks down in ``large'' economies: in private goods economies the incentive to misrepresent dwindles as the economy increases, but this is not the case for public goods economies (J. Roberts [1976], J. Roberts and Postlewaite [1976]).

\section*{Economies with Production}

Private goods economies with production provided some of the original stimulus for the study of alternative mechanisms and incentive-compatibility games. Thus, for instance, in a socialist economy of the Lange--Lerner type, if the environment is ``classical,'' each manager of a production unit is supposed to choose a profit-maximizing input–output vector \textit{treating prices parametrically}.\footnote{More generally, when the production function is not assumed concave, the requirement is that the input–output vector yield the equality of output prices and marginal costs, again with all prices treated parametrically.} (The result is a competitive equilibrium.) But if, for instance, the firm has a monopoly of its product, higher profits could be attained by exploiting the monopoly power---that is, by \textit{not} treating prices parametrically. Now suppose the manager's reward is an increasing function of profits. This creates an incentive to misrepresent the firm's production function in such a way that a level of output that maximizes monopoly profits for the true production function is presented as the competitive profit-maximizing output for a false production function (Hurwicz [1973]). Thus, in an economy with a finite number of firms, the competitive mechanism is not incentive compatible if the decision maker's rewards increase with profits and informational decentralization prevails.\footnote{It is of particular interest to note that imposing rules of behavior (such as profit maximization) where actions are prescribed as functions of the agents' characteristics, and with the center acting on the assumption that these rules are honestly obeyed, is equivalent to operating a natural direct-revelation game!} But, again, it is possible to design mechanisms implementing a wide class of performance correspondences (Hurwicz, Maskin, and Postlewaite [1980]). These mechanisms use (as strategies) the profiles $(Y^1, \dots, Y^n)$ of production possibility sets, similarly to the use of preference profiles in Maskin [1977]. This problem is somewhat analogous to that arising in economies in which endowments may be destroyed by agents. The analysis of the latter problem was pioneered by Postlewaite [1979].

It may be noted that the presence of production creates difficulties in reconciling efficiency with certain notions of fairness (e.g., in the sense of being envy free). Much of Pazner's work was devoted to this problem.

\medskip

By now the literature dealing with economic mechanisms is huge. The present notes do no more than pursue a few topics that seem of particular importance in the earlier stages of development of this area. Most workers in the field are far from satisfied with the accomplishments to date. There are differences of opinion as to the appropriateness of different game-theoretic models. There is a need for better integration of incentive and informational aspects of our models. Also, the issue of combining fairness with efficiency is still before us. Elisha Pazner's pioneering insights in this area will continue to illuminate the path for future explorers.

\vspace{1.5cm}

% --- Bibliografía Final ---
\begin{center}
    {\large\scshape\color{MidnightBlue} References}
    \vspace{0.8cm}
\end{center}

\bibentry{Arrow, K. J. [1951]. ``An Extension of the Basic Theorems of Welfare Economics,'' in \textit{Proceedings of the Second Berkeley Symposium on Mathematical Statistics and Probability} edited by G. Neyman, pp. 507--532. Berkeley: University of California Press.}

\bibentry{Arrow, K. J. [1951, 1963]. \textit{Social Choice and Individual Values}. New York: Wiley, 1951. New Haven: Yale University Press, 1963.}

\bibentry{Arrow, K. J. [1977]. ``The Property Rights Doctrine and Demand Revelation under Incomplete Information.'' IMSSS Technical Report no. 243, Stanford University.}

\bibentry{Arrow, K. J., and F. Hahn [1971]. \textit{General Competitive Analysis}. San Francisco: Holden-Day.}

\bibentry{Bergson, A. [1938]. ``A Reformulation of Certain Aspects of Welfare Economics.'' \textit{Quarterly Journal of Economics}, 52:310--334.}

\bibentry{Clarke, E. H. [1971]. ``Multipart Pricing of Public Goods.'' \textit{Public Choice}, 2:19--33.}

\bibentry{Dasgupta, P., P. Hammond, and E. Maskin [1979]. ``The Implementation of Social Choice Rules: Some General Results on Incentive Compatibility.'' \textit{Review of Economic Studies}, 46:181--216.}

\bibentry{d'Aspremont, C., and L.-A. Gérard-Varet [1979a]. ``On Bayesian Incentive Compatible Mechanisms,'' in \textit{Aggregation and Revelation of Preferences}, edited by J.-J. Laffont, Ch. 22, pp. 269--288. Amsterdam: North Holland.}

\bibentry{d'Aspremont, C., and L.-A. Gérard-Varet [1979b]. ``Incentives and Incomplete Information,'' \textit{Journal of Public Economics}, 11:25--45.}

\bibentry{Debreu, G. [1959]. \textit{Theory of Value}. New York: Wiley.}

\bibentry{Drèze, J., and D. de la Vallée Poussin [1969]. ``A Tâtonnement Process for Guiding and Financing an Efficient Production for Public Goods,'' CORE paper, 1969. (\textit{RES}, 38 [1971], pp. 133--150.)}

\bibentry{Foley, D. K. [1970]. ``Lindahl's Solution and the Core of an Economy with Public Goods.'' \textit{Econometrica}, 38:66--72.}

\bibentry{Green, J., and J.-J. Laffont [1977]. ``Characterization of Satisfactory Mechanisms for the Revelation of Preferences for Public Goods.'' \textit{Econometrica}, 45: 427--438.}

\bibentry{Green, J., and J.-J. Laffont [1979]. \textit{Incentives in Public Decision-Making}. Amsterdam: North Holland.}

\bibentry{Groves, T. [1970]. ``The Allocation of Resources under Uncertainty.'' Unpublished Ph.D. dissertation, University of California, Berkeley.}

\bibentry{Groves, T., and J. Ledyard [1977]. ``Optimal Allocation of Public Goods: A Solution to the `Free Rider' Problem.'' \textit{Econometrica}, 45(4):783--811.}

\bibentry{Groves, T., and J. Ledyard [1980]. ``The Existence of Efficient and Incentive Compatible Equilibria with Public Goods.'' \textit{Econometrica}, 48:1487--1506.}

\bibentry{Groves, T., and M. Loeb [1975]. ``Incentives and Public Inputs.'' \textit{Journal of Public Economies}, 4(3):211--226.}

\bibentry{Harsanyi, J. [1967--68]. ``Games with Incomplete Information Played by Bayesian Players.'' \textit{Management Science} 14:159--82, 320--34, 486--502.}

\bibentry{von Hayek, F.A. [1945]. ``The Use of Knowledge in Society.'' \textit{American Economic Review}, 35:519--530.}

\bibentry{Holmström, B. [1979]. ``Groves Scheme on Restricted Domain.'' \textit{Econometrica}, 47: 1137--1144.}


\bibentry{Hotelling, H. [1938]. ``The General Welfare in Relation to Problems of Taxation and of Railway and Utility Rates.'' \textit{Econometrica}, 6:242--269.}

\bibentry{Hurwicz, L. [1972]. ``On Informationally Decentralized Systems,'' in \textit{Decision and Organization (Volume in Honor of J. Marschak)}, edited by R. Radner and C. B. McGuire, pp. 297--336. Amsterdam: North Holland.}

\bibentry{Hurwicz, L. [1973]. ``The Design of Mechanisms for Resource Allocation.'' \textit{American Economic Review}, 63:1--30, esp. pp. 25--26.}

\bibentry{Hurwicz, L. [1977]. ``On the Dimensional Requirements of Informationally Decentralized Pareto-Satisfactory Processes,'' in \textit{Studies in Resource Allocation Processes}, edited by K. J. Arrow and L. Hurwicz, pp. 413--424. Cambridge: Cambridge University Press.}

\bibentry{Hurwicz, L. [1979a]. ``Outcome Functions Yielding Walrasian and Lindahl Allocations at Nash Equilibrium Points.'' \textit{Review of Economic Studies}, 46(2):217--225.}

\bibentry{Hurwicz, L. [1979b]. ``On Allocations Attainable through Nash Equilibria.'' \textit{Journal of Economic Theory}, 21, (1):140--165; also in: \textit{Aggregation and Revelation of Preferences}, edited by J.-J. Laffont, Ch. 22, pp. 397--419. Amsterdam: North Holland.}

\bibentry{Hurwicz, L. [1979c]. ``Balanced Outcome Functions Yielding Walrasian and Lindahl Allocations at Nash Equilibrium Points for Two or More Agents,'' in \textit{General Equilibrium, Growth, and Trade}, edited by R. Green and J. A. Scheinkman. New York: Academic Press.}

\bibentry{Hurwicz, L., E. Maskin, and A. Postlewaite [1980]. ``Feasible Implementation of Social Choice Correspondences by Nash Equilibria.'' Mimeo. [Abbreviated as HMP.]}

\bibentry{Hurwicz, L., and M. Walker [1983]. ``On the Generic Non-Optimality of Dominant-Strategy Mechanisms with an Application to Pure Exchange.'' Research Paper No. 250, circulated by the Economic Research Bureau, State University of New York at Stonybrook.}

\bibentry{Johansen, L. [1963]. ``Some Notes on the Lindahl Theory of Determination of Public Expenditure,'' \textit{International Economic Review}, 4:346--358.}

\bibentry{Jordan, J. S. [1982]. ``The Competitive Allocation Process Is Informationally Efficient Uniquely.'' \textit{Journal of Economic Theory}, 28:1--18.}

\bibentry{Koopmans, T. C. [1957]. ``Allocation of Resources and the Price System,'' in \textit{Three Essays on the State of Economic Science}. New York: McGraw-Hill.}

\bibentry{Lange, O., and F. M. Taylor [1938]. \textit{On the Economic Theory of Socialism}. Minneapolis: University of Minnesota Press.}

\bibentry{Lange, O. [1942]. ``The Foundations of Welfare Economics.'' \textit{Econometrica}, 10:215--228.}

\bibentry{Ledyard, J. [1977]. ``Incentive Compatible Behavior in Core-Selecting Organizations.'' \textit{Econometrica}, 45(7):1607--1621.}

\bibentry{Ledyard, J., and J. Roberts [1974]. ``On the Incentive Problems with Public Goods.'' Discussion Paper No. 116, CMSEMS, Northwestern University, Evanston, Ill.}

\bibentry{Lerner, A. P. [1944]. \textit{Economics of Control}. New York: Macmillan.}

\bibentry{Lindahl, E. [1919]. ``Just Taxation---A Positive Solution,'' in \textit{Classics in the Theory of Public Finance}, edited by R. A. Musgrave and A. T. Peacock, London: Macmillan, 1964, pp. 168--176. (Also, New York: St. Martin's Press, 1964.)}

\bibentry{Malinvaud, E. [1969]. ``Procédures pour la determinación d'un programme de consommation collective.'' Paper presented at the Brussels meeting of the Econometric Society, September 1969. Mimeo.}

\bibentry{Maskin, E. [1977]. ``Nash Equilibrium and Welfare Optimality.'' Working paper, October 1977, Massachusetts Institute of Technology; scheduled to appear in \textit{Mathematics of Operations Research}.}

\bibentry{Milleron, J.-C. [1972]. ``Theory of Value with Public Goods: A Survey Article.'' \textit{Journal of Economic Theory}, 5:419--477.}

\bibentry{Miura, Rei [1982]. ``Counter-Example and Revised Outcome Function Yielding the Nash-Lindahl Equivalence for Two or More Agents.'' Keio University, Tokyo. Privately circulated paper.}

\bibentry{Mount, K., and S. Reiter [1974]. ``The Informational Size of Message Spaces.'' \textit{Journal of Economic Theory}, 8(3):161--192.}

\bibentry{Osana, H. [1978]. ``On the Informational Size of Message Spaces for Resource Allocation Processes.'' \textit{Journal of Economic Theory}, 17:66--78.}

\bibentry{Postlewaite, A. [1979]. ``Manipulation via Endowments.'' \textit{Review of Economic Studies}, 46:255--262.}

\bibentry{Postlewaite, A., and D. Wettstein [1983]. ``Implementing Constrained Walrasian Equilibria Continuously.'' CARESS Working Paper \#83-24, University of Pennsylvania.}

\bibentry{Reichelstein, S. [1982]. ``On the Informational Requirements for the Implementation of Social Choice Rules.'' Handout for the Decentralized Conference, May.}

\bibentry{Reichelstein, S. [1984]. ``Smooth versus Discontinuous Mechanisms.'' February, privately circulated paper.}

\bibentry{Reiter, S. [1977]. ``Information and Performance in the (New)$^2$ Welfare Economics.'' \textit{American Economic Review}, 67(1):226--234.}

\bibentry{Roberts, J. [1976]. ``The Incentives for Correct Revelation of Preferences and the Number of Consumers.'' \textit{Journal of Public Economics}, 6:359--374.}

\bibentry{Roberts, J., and A. Postlewaite [1976]. ``The Incentives for Price-Taking Behavior in Large Economies.'' \textit{Econometrica}, 44:115--128.}

\bibentry{Samuelson, P. [1954]. ``The Pure Theory of Public Expenditure.'' \textit{The Review of Economics and Statistics}, 36:387--389.}

\bibentry{Samuelson, P. [1955]. ``Diagrammatic Exposition of a Theory of Public Expenditure.'' \textit{The Review of Economics and Statistics} 37:350--356.}

\bibentry{Sato, F. [1981]. ``On the Informational Size of Message Spaces for Resource Allocation Processes in Economies with Public Goods,'' \textit{Journal of Economic Theory}, 24(1):48--69.}

\bibentry{Satterthwaite, M. [1976]. ``Straightforward Allocation Mechanisms.'' Discussion Paper no. 253, CMSEMS, Northwestern University, Evanston, Ill. October.}

\bibentry{Satterthwaite, M., and H. Sonnenschein [1981]. ``Strategy Proof Allocation Mechanisms at Differentiable Points.'' \textit{Review of Economic Studies}, 48:587--597.}

\bibentry{Schmeidler, D. [1976]. ``A Remark on a Game Theoretic Interpretation of Walras Equilibria.'' University of Minnesota, Minneapolis. Mimeo.}

\bibentry{Schmeidler, D. [1980]. ``Walrasian Analysis via Strategic Outcome Functions.'' \textit{Econometrica}, 48:1585--1594.}

\bibentry{Schmeidler, D. [1982]. ``A Condition Guaranteeing that the Nash Allocation Is Walrasian.'' \textit{Journal of Economic Theory}, 28:376--378.}


\bibentry{Sen, A. K. [1970]. \textit{Collective Choice and Social Welfare}. San Francisco: Holden-Day.}

\bibentry{Thomson, W. [1979]. ``Comment on L. Hurwicz: On Allocations Attainable through Nash Equilibria,'' in \textit{Aggregation and Revelation of Preferences}, edited by J.-J. Laffont, pp. 420--431. Amsterdam: North Holland.}

\bibentry{Walker, M. [1980]. ``On the Nonexistence of a Dominant Strategy Mechanism for Making Optimal Public Decisions.'' \textit{Econometrica}, 48(6):1521--1540.}

\bibentry{Walker, M. [1981]. ``A Simple Incentive Compatible Scheme for Attaining Lindahl Allocations.'' \textit{Econometrica}, 49:65--73.}


\end{document}
